\documentclass[11pt]{article}
\usepackage[spanish]{babel}
\usepackage{amsmath}
\usepackage{amssymb}
\usepackage{graphicx}
\graphicspath{./images/}
\usepackage[utf8]{inputenc}
\usepackage[T1]{fontenc}
\usepackage[margin=1in]{geometry}
\usepackage{fancyhdr}
\usepackage{hyperref}

\hypersetup{
    pdftitle={Comunicación Efectiva en Contextos Académicos},
    pdfauthor={Felipe Colli Olea},
    colorlinks=true,
    linkcolor=black,
    urlcolor=cyan
}

\pagestyle{fancy}
\fancyhf{}
\fancyfoot[C]{\footnotesize Felipe Colli, Todos los Derechos Reservados}
\fancyfoot[R]{\thepage}
\renewcommand{\headrulewidth}{0pt}

% Corrección de tildes y mayúsculas según las instrucciones
\title{Comunicación Efectiva en Contextos Académicos}
\author{Felipe Colli Olea}
\date{\today}

\begin{document}

\maketitle
\thispagestyle{fancy}

% Opcional: El índice puede ser un poco excesivo para un documento breve,
% pero lo mantengo ya que lo incluiste.
\tableofcontents

\newpage

\section{Informes de Laboratorio}

El informe de laboratorio es un género discursivo cuyo propósito es mostrar y evidenciar los resultados obtenidos de una experiencia experimental, analizarlos y contrastarlos con los resultados esperados teóricamente.

La estructura general de un informe de laboratorio se compone de los siguientes apartados:
\begin{enumerate}
	\item \textbf{Resumen ejecutivo:} Síntesis de todo el texto (tema, antecedentes, resultados cualitativos, análisis y conclusiones principales).
	\item \textbf{Introducción:} Contextualización para el lector sobre la experiencia realizada.
	\item \textbf{Objetivos:} Se dividen en generales (qué resultados se deben obtener) y específicos (hitos para cumplir el objetivo general).
	\item \textbf{Antecedentes:} Base conceptual que fundamenta el análisis. Se dividen en teóricos (principios físicos o matemáticos) y prácticos (variables a medir y equipos utilizados).
	\item \textbf{Metodología:} Descripción de los pasos realizados y precauciones en la manipulación de los equipos.
	\item \textbf{Memoria de cálculo:} Enumeración de las fórmulas utilizadas y glosario de términos con sus unidades.
	\item \textbf{Resultados:} Presentación de los datos medidos y resultados de cálculos, apoyados visualmente por tablas y gráficos.
	\item \textbf{Análisis o Discusión de resultados:} Interpretación basada en la teoría, evaluación de la metodología y posibles causas de errores.
	\item \textbf{Conclusiones:} Síntesis de los hallazgos y evaluación del cumplimiento de los objetivos.
	\item \textbf{Referencias:} Detalle de las fuentes bibliográficas utilizadas.
\end{enumerate}

Para una escritura efectiva, se recomienda redactar en tercera persona (lenguaje impersonal), mantener un registro formal, evitar valoraciones subjetivas y ser explícito al describir fenómenos. Además, los párrafos deben tener una estructura clara compuesta por una oración principal, oraciones de desarrollo y una oración de cierre.

\newpage
\section{Informes de Proyecto}

Un informe de proyecto es un documento técnico utilizado en la ingeniería y las ciencias para proponer la ejecución de un proceso o la innovación en actividades de desarrollo. Su objetivo principal es dar solución a una problemática detectada, justificándola debidamente para lograr la aprobación técnica y de presupuesto por parte de una entidad.

La información en este tipo de documento se organiza de la siguiente manera:
\begin{enumerate}
	\item \textbf{Resumen del proyecto:} Introducción al lector sobre el tema y los componentes principales de la propuesta.
	\item \textbf{Datos de la entidad solicitante:} Identificación de la organización o miembros que proponen el proyecto.
	\item \textbf{Descripción del proyecto:} Aporta datos sobre el contexto, el problema específico a solucionar, la relevancia de la propuesta, los procedimientos a llevar a cabo y los resultados esperados.
	\item \textbf{Viabilidad e impacto del proyecto:} Detección de las posibles dificultades en el desarrollo y de las ganancias o impactos positivos que traerá su ejecución.
	\item \textbf{Presupuesto:} Detalle exhaustivo de los recursos económicos solicitados.
	\item \textbf{Referencias bibliográficas:} Respaldos teóricos que sustentan la propuesta.
\end{enumerate}

En la redacción de proyectos es fundamental cuidar la \textbf{coherencia} (organización lógica de las ideas, como el patrón Problema-Solución) y la \textbf{cohesión} (relaciones gramaticales para evitar la repetición y facilitar la lectura). Algunos mecanismos útiles de cohesión incluyen la sustitución por pronombres, el uso de sinónimos o hiperónimos, y la elipsis (omisión de un elemento ya mencionado que queda implícito).

\newpage
\section{Artículos de Divulgación Científica}

El conocimiento científico suele circular en contextos sumamente especializados. El artículo de divulgación científica nace con el desafío de \textit{recontextualizar} dicho conocimiento para orientarlo a una audiencia general, menos especializada. Su propósito es acercar el saber tecnológico e investigativo a la sociedad, explicando el conocimiento de manera comprensible.

La estructura de un artículo de divulgación se organiza en:
\begin{itemize}
	\item \textbf{Resumen del artículo:} Resume los antecedentes del tema y anticipa la contribución que este tiene para la sociedad.
	\item \textbf{Desarrollo o cuerpo del texto:} Contextualiza los nuevos hallazgos, define y explica conceptos utilizando un lenguaje reformulado para el público general, y valora la importancia del avance para el conocimiento.
	\item \textbf{Información complementaria:} Profundiza en perspectivas históricas o entrega enlaces de referencia para quienes deseen indagar más en el tema.
\end{itemize}

El éxito de este documento radica en la \textbf{recontextualización de la práctica social}. Para que un texto experto pase a ser divulgativo, el escritor debe transformar los siguientes elementos:
\begin{itemize}
	\item \textbf{Participantes:} El receptor pasa de ser una audiencia experta a un público de carácter más general.
	\item \textbf{Acciones:} Los propósitos cambian de informar entre pares académicos a explicar detalladamente y divulgar.
	\item \textbf{Estilos de presentación:} Se reemplaza el uso excesivo de fórmulas y jerga técnica por un lenguaje amigable y contextualizado.
	\item \textbf{Recursos:} Se aumenta el uso de recursos visuales (imágenes, esquemas) y verbales (explicaciones detalladas, analogías) para facilitar la comprensión de conceptos complejos.
\end{itemize}

\end{document}
